\section{Mid-Term Examination --- Set D}
\label{sec:midterm_set_d}

\LPUPaperHeader{Mid-Term Examination (Objective MCQ)}{D}{90 Minutes}{30 Marks}{CO1 (Unit 1: Matrices), CO2 (Unit 2: ODEs \& Calculus), CO3 (Unit 3: Infinite Series)}

\begin{center}
\sffamily\textbf{\color{AlertRed}Instructions: All 30 questions are compulsory. Each question carries 1 mark. A penalty of 0.25 marks is deducted for each incorrect response.}
\end{center}

\vspace{3mm}

% ------------------------------------------------------------------------------
% 30 MCQS
% ------------------------------------------------------------------------------

\mcqitem{1}{The eigenvalues of a triangular matrix are:}{All real and equal}{The diagonal elements themselves}{Purely imaginary}{Reciprocals of the diagonal elements}{1}

\mcqitem{2}{If $A = \begin{bmatrix} 1 & 2 \\ 0 & 3 \end{bmatrix}$, then the eigenvalues of $A^{-1}$ are:}{$1, 2$}{$1, 1/3$}{$1/2, 1/3$}{$-1, -3$}{1}

\mcqitem{3}{The value of the limit $\lim_{x\to 0} \frac{\tan x - x}{x^3}$ is:}{$1/3$}{$1/6$}{$0$}{$\infty$}{1}

\mcqitem{4}{The series $\sum_{n=1}^\infty \frac{n^2}{2^n}$ is:}{Divergent}{Convergent by Ratio Test}{Oscillatory}{Conditionally Convergent}{1}

\mcqitem{5}{The complementary function of $(D^3 - 6D^2 + 11D - 6)y = 0$ is:}{$c_1 e^x + c_2 e^{2x} + c_3 e^{3x}$}{$(c_1 + c_2 x + c_3 x^2)e^x$}{$c_1 e^{-x} + c_2 e^{-2x} + c_3 e^{-3x}$}{$c_1 \cos x + c_2 \cos 2x + c_3 \cos 3x$}{1}

\mcqitem{6}{If $A$ is a skew-Hermitian matrix, then $iA$ is:}{Skew-Hermitian}{Hermitian}{Orthogonal}{Unitary}{1}

\mcqitem{7}{The value of the definite integral $\int_0^{\pi/2} \sin^3 x \, dx$ is:}{$1/3$}{$2/3$}{$\pi/4$}{$\pi/6$}{1}

\mcqitem{8}{The particular integral of $(D^2 + 16)y = e^{3x}$ is:}{$\frac{e^{3x}}{25}$}{$\frac{e^{3x}}{7}$}{$\frac{e^{3x}}{9}$}{$\frac{x e^{3x}}{25}$}{1}

\mcqitem{9}{The series $\sum_{n=1}^\infty \frac{1}{n^3 + n}$ converges by comparison with:}{$\sum \frac{1}{n}$}{$\sum \frac{1}{n^2}$}{$\sum \frac{1}{n^3}$}{$\sum \frac{1}{n!}$}{1}

\mcqitem{10}{If the rank of a $3\times 3$ matrix $A$ is 2, then $\det(A)$ is:}{$1$}{$-1$}{$0$}{Any non-zero real number}{1}

\mcqitem{11}{The Maclaurin series of $\cos x$ contains only:}{Odd powers of $x$}{Even powers of $x$}{Negative powers of $x$}{Fractional powers of $x$}{1}

\mcqitem{12}{The particular integral of $(D^2 - 4)y = \cosh 2x$ is:}{$\frac{x}{4}\sinh 2x$}{$\frac{x}{2}\sinh 2x$}{$\frac{1}{4}\cosh 2x$}{$\frac{x^2}{8}\cosh 2x$}{1}

\mcqitem{13}{The radius of curvature of a straight line $y = mx + c$ is:}{$0$}{$1$}{$\infty$}{$m$}{1}

\mcqitem{14}{A matrix $A$ is unitary if:}{$A^T A = I$}{$A^\dagger A = I$}{$A^2 = I$}{$A = A^\dagger$}{1}

\mcqitem{15}{The alternating harmonic series $\sum_{n=1}^\infty \frac{(-1)^{n-1}}{n}$ has the exact sum:}{$\ln 2$}{$\ln 3$}{$1$}{$0$}{1}

\mcqitem{16}{The general solution of $y'' + 9y = 0$ is:}{$c_1 e^{3x} + c_2 e^{-3x}$}{$c_1 \cos 3x + c_2 \sin 3x$}{$(c_1 + c_2 x)e^{3x}$}{$c_1 \cosh 3x + c_2 \sinh 3x$}{1}

\mcqitem{17}{The rank of the zero matrix $O_{m\times n}$ is:}{$1$}{$0$}{$m$}{$n$}{1}

\mcqitem{18}{The value of $\Gamma(1/2)$ is:}{$\sqrt{\pi}$}{$\pi$}{$\frac{\pi}{2}$}{$1$}{1}

\mcqitem{19}{The particular integral of $(D^2 + 1)y = x^2$ is:}{$x^2 - 2$}{$x^2 + 2$}{$x^2$}{$\frac{x^3}{3}$}{1}

\mcqitem{20}{If the characteristic roots of a matrix are $1, 2, 3$, the determinant of the matrix is:}{$6$}{$5$}{$0$}{$14$}{1}

\mcqitem{21}{The series $\sum_{n=1}^\infty \frac{1}{n(n+1)}$ has the sum:}{$1$}{$1/2$}{$2$}{$\infty$}{1}

\mcqitem{22}{The order of the differential equation $\left(\frac{d^3 y}{dx^3}\right)^2 + \left(\frac{d^2 y}{dx^2}\right)^5 + y = 0$ is:}{$2$}{$3$}{$5$}{$6$}{1}

\mcqitem{23}{The trace of an orthogonal matrix of order 3 can never exceed:}{$1$}{$2$}{$3$}{$9$}{1}

\mcqitem{24}{The convergence of $\sum_{n=1}^\infty \frac{1}{n!}$ is established by:}{Ratio Test}{p-Series Test}{Raabe's Test}{Integral Test}{1}

\mcqitem{25}{The solution of $(D - 3)^3 y = 0$ is:}{$(c_1 + c_2 x + c_3 x^2)e^{3x}$}{$c_1 e^{3x} + c_2 e^{-3x} + c_3$}{$(c_1 + c_2 x)e^{3x}$}{$c_1 e^x + c_2 e^{2x} + c_3 e^{3x}$}{1}

\mcqitem{26}{If $A$ is a real symmetric matrix, its eigenvectors corresponding to distinct eigenvalues are:}{Parallel}{Orthogonal}{Linearly Dependent}{Zero}{1}

\mcqitem{27}{The value of $\lim_{x\to 0} \frac{e^x - 1}{x}$ is:}{$0$}{$1$}{$e$}{$\infty$}{1}

\mcqitem{28}{The series $\sum_{n=1}^\infty \frac{(-1)^n}{n^3}$ converges:}{Conditionally only}{Absolutely}{Does not converge}{Oscillates}{1}

\mcqitem{29}{The general solution of $\frac{dy}{dx} + y = 0$ is:}{$y = c e^{-x}$}{$y = c e^x$}{$y = -x + c$}{$y = c \ln x$}{1}

\mcqitem{30}{The matrix equation $AX = 0$ always has:}{No solution}{At least the trivial solution $X = 0$}{Only non-trivial solutions}{Infinitely many solutions only}{1}

\vspace{5mm}
\hrule
\vspace{5mm}

% ------------------------------------------------------------------------------
% ANSWER KEY & EXPLANATIONS
% ------------------------------------------------------------------------------
\subsection*{Answer Key (Mid-Term Set D)}

\begin{center}
\begin{tabular}{|c|c||c|c||c|c||c|c||c|c||c|c|}
\hline
\textbf{Q} & \textbf{Ans} & \textbf{Q} & \textbf{Ans} & \textbf{Q} & \textbf{Ans} & \textbf{Q} & \textbf{Ans} & \textbf{Q} & \textbf{Ans} & \textbf{Q} & \textbf{Ans} \\
\hline
1 & \textbf{B} & 6 & \textbf{B} & 11 & \textbf{B} & 16 & \textbf{B} & 21 & \textbf{A} & 26 & \textbf{B} \\
2 & \textbf{B} & 7 & \textbf{B} & 12 & \textbf{A} & 17 & \textbf{B} & 22 & \textbf{B} & 27 & \textbf{B} \\
3 & \textbf{A} & 8 & \textbf{A} & 13 & \textbf{C} & 18 & \textbf{A} & 23 & \textbf{C} & 28 & \textbf{B} \\
4 & \textbf{B} & 9 & \textbf{C} & 14 & \textbf{B} & 19 & \textbf{A} & 24 & \textbf{A} & 29 & \textbf{A} \\
5 & \textbf{A} & 10 & \textbf{C} & 15 & \textbf{A} & 20 & \textbf{A} & 25 & \textbf{A} & 30 & \textbf{B} \\
\hline
\end{tabular}
\end{center}

\vspace{3mm}
\subsection*{Selected Detailed Mathematical Explanations}

\begin{solutionbox}[Explanations for Critical Questions]
\textbf{Q.2 (Ans: B)} Eigenvalues of upper triangular $A$ are $1$ and $3$. Eigenvalues of $A^{-1}$ are reciprocals: $1/1 = 1$ and $1/3$.\\[2mm]
\textbf{Q.3 (Ans: A)} Using series expansion: $\tan x = x + x^3/3 + \dots \implies \frac{\tan x - x}{x^3} = 1/3$.\\[2mm]
\textbf{Q.4 (Ans: B)} $\lim_{n\to\infty} \frac{u_{n+1}}{u_n} = \lim \frac{(n+1)^2}{2^{n+1}} \frac{2^n}{n^2} = \frac{1}{2} < 1 \implies$ Converges.\\[2mm]
\textbf{Q.6 (Ans: B)} If $A^\dagger = -A$, then $(iA)^\dagger = -i A^\dagger = -i(-A) = iA$, which is Hermitian.\\[2mm]
\textbf{Q.12 (Ans: A)} $\cosh 2x = \frac{e^{2x} + e^{-2x}}{2}$. For $\frac{1}{D^2-4}e^{2x} = \frac{x}{2(2)}e^{2x} = \frac{x}{4}e^{2x}$. Similarly for $e^{-2x}$ gives $-\frac{x}{4}e^{-2x}$. Adding gives $\frac{x}{4}\sinh 2x$.\\[2mm]
\textbf{Q.18 (Ans: A)} Standard Gamma value: $\Gamma(1/2) = \int_0^\infty t^{-1/2}e^{-t}dt = \sqrt{\pi}$.\\[2mm]
\textbf{Q.19 (Ans: A)} $\frac{1}{1+D^2}x^2 = (1 - D^2 + D^4)x^2 = x^2 - 2$.\\[2mm]
\textbf{Q.21 (Ans: A)} Telescoping sum: $\sum_{n=1}^\infty \left(\frac{1}{n} - \frac{1}{n+1}\right) = (1 - 1/2) + (1/2 - 1/3) + \dots = 1$.
\end{solutionbox}

\newpage
